% META: {"id": "TSPE-COMBINATOIRE-COURS-01", "chapitre": "TSPE-COMBINATOIRE", "type_objet": "cours", "capacites_codes": ["C1"], "sources_inspiration": [], "mode_creation": "ex_nihilo", "status": "approved", "capacites": ["TSPE-COMBINATOIRE-C1"], "statut": "structure"}

\section{Représentations pour dénombrer}

% ============================================================
% STRATE 1 — Outils de representation
% ============================================================

\definition[1]{
\textbf{Dénombrer} un ensemble fini, c'est determiner son nombre d'elements (son \textbf{cardinal}), note $\mathrm{Card}(E)$ ou $|E|$.
}

\propriete[Outils de représentation]{
\begin{itemize}
  \item \textbf{Arbre} : represente des choix successifs (chaque branche = une possibilite a une etape).
  \item \textbf{Tableau} (double entree) : adapte pour deux choix independants.
  \item \textbf{Diagramme (patates, Venn)} : adapte pour representer des unions, intersections, complementaires d'ensembles.
  \item \textbf{Liste ou ensemble} : enumeration directe, adaptee aux petits cas.
\end{itemize}
}

\exemple{
Un menu propose $3$ entrees et $2$ desserts. Combien de menus (entree, dessert) different differents peut-on composer ?

Un arbre a $3$ branches au premier niveau (entrees), chacune se subdivisant en $2$ branches (desserts) : on obtient $3 \times 2 = 6$ chemins, donc $6$ menus possibles.

% BEGIN-VERIFY
% assert 3 * 2 == 6
% END-VERIFY
}

\margeAppui{
Le choix de la representation depend de la structure du probleme : arbre pour des etapes successives, tableau pour deux criteres croises, diagramme pour des categories qui se recoupent.
}

% ============================================================
% STRATE 2 — Erreurs frequentes
% ============================================================

\erreurFrequente{
\textbf{Oublier de verifier si l'ordre compte et si les repetitions sont autorisees.} Avant tout denombrement, il faut se demander : "Choisit-on avec ou sans ordre ?" et "Peut-on choisir plusieurs fois le meme element ?" — ces deux questions determinent la formule a utiliser (chapitre suivant).
}
